\begin{theorem}
    \label{teo:polinomios_weyl}
    Sea $p(x)=c_m x^m + c_{m-1} x^{m-1} + \dots + c_1 x + c_0$ un polinomio con coeficientes 
    reales. Si al menos uno de $c_j \notin \mathbb{Q}$ con $j>0$, entonces la sucesión $(p(n))$ 
    es equidistribuida módulo $1$.
\end{theorem}

\begin{proof}
    Lo primero que haremos será comprobar el caso especial en el que $c_1 \notin \mathbb{Q}$ 
    y los demás coeficientes son racionales.
    Sea
    \[
        p(x) = P(x) + c_1x + c_0
    \]
    con $P(x) = c_m x^m + c_{m-1} x^{m-1} + \dots + c_2 x^2$.

    Como $c_2,\dots,c_m \in \mathbb{Q}$ se pueden escribir de la forma $c_j = p_j/q_j$ con 
    $p_j,q_j \in \mathbb{Z}$ y $q_j > 0$ para $j=~2,\dots,m$. 
    Sea $D = \operatorname{mcm}(q_2,\dots,q_m)$ entonces
    \[
        c_j = \frac{a_j}{D} \quad \text{con } a_j \in \mathbb{Z} \text{ para cada } j=2,\dots,m
    \]

    Sea un $n\geq1$, podemos expresarlo como $n=Dl+d$ con $l \geq 0$ y $1 \leq d \leq D$.
    Si comparamos $P(Dl+d)$ con $P(d)$ se tiene que
    \[
        P(Dl+d)-P(d)= \sum_{j=2}^m c_j ((Dl+d)^j - d^j)
    \]
    Para cada $j=2,...,m$, desarrollando el binomio,
    \[
        (Dl+d)^j - d^j = \sum_{i=1}^j \binom{j}{i} d^{j-i} (Dl)^i
    \]
    Por lo tanto,
    \begin{align*}
        P(Dl+d)-P(d) 
        &= \sum_{j=2}^m \sum_{i=1}^j c_j \binom{j}{i} d^{j-i} (Dl)^i \\
        &= \sum_{j=2}^m \sum_{i=1}^j \frac{a_j}{D} \binom{j}{i} d^{j-i} (Dl)^i \\
        &= \sum_{j=2}^m \sum_{i=1}^j a_j \binom{j}{i} d^{j-i} D^{i-1} l^i \in \mathbb{Z}
    \end{align*}
    Es decir, $P(Dl+d)$ y $P(d)$ difieren en un número entero, por lo que
    \[
        \partfrac{P(Dl+d)} = \partfrac{P(d)} 
        \quad \text{y en consecuencia} \quad
        e^{2\pi i P(Dl+d)} = e^{2\pi i P(d)}
    \]

    Comprobaremos que para el caso $c_1 \notin \mathbb{Q}$ la sucesión 
    $(p(n))$ es equidistribuida módulo $1$ 

    Por el criterio de Weyl, sea $k \in \mathbb{Z}\setminus\{0\}$, sea $N \geq 1$, 
    entonces $N=DL+R$ con $L \geq 0$ y $1 \leq R \leq D$. Para cada $n=1,2,\dots,DL$
    lo expresamos como $n= Dl + d$ con $0 \leq l \leq L-1$ y $1 \leq d \leq D$. Entonces
    \begin{align*}
        \left|\frac{1}{N}\sum_{n=1}^N e^{2\pi i k p(n)} \right|
        &= \left|\frac{1}{N}\sum_{d=1}^D \sum_{l=0}^{L-1} e^{2\pi i k(P(Dl+d) + c_1 Dl + c_0)} + \frac{1}{N}\sum_{n=DL+1}^N e^{2\pi i k p(n)} \right| \\
        &\leq \left|\frac{1}{N}\sum_{d=1}^D e^{2\pi i k(P(d) + c_0)} \sum_{l=0}^{L-1} e^{2\pi i k (c_1 Dl)}\right| + \frac{1}{N}\sum_{n=DL+1}^N \left|e^{2\pi i k p(n)}\right| \\
        &\leq \left|\frac{D}{N} \sum_{l=0}^{L-1} e^{2\pi i k (c_1 Dl)}\right| + \frac{D}{N}
    \end{align*}
    Como $N=DL+R$ con $1 \leq R \leq D$, entonces $\frac{D}{N} \leq \frac{D}{DL} = \frac{1}{L}$, por lo que
    \[
        \left|\frac{1}{N}\sum_{n=1}^N e^{2\pi i k p(n)} \right| \leq 
        \left|\frac{1}{L} \sum_{l=0}^{L-1} e^{2\pi i k (c_1 Dl)}\right| + \frac{1}{L}
    \]

    Dado que $c_1 \notin \mathbb{Q}$, entonces $(c_1 Dl)$ es equidistribuida módulo $1$.
    Al tomar el límite cuando $N \to \infty$, y por lo tanto $L \to \infty$, obtenemos
    \[
        \lim_{N \to \infty} \left|\frac{1}{N}\sum_{n=1}^N e^{2\pi i k p(n)} \right| \leq 
        \lim_{L \to \infty} \left|\frac{1}{L} \sum_{l=0}^{L-1} e^{2\pi i k (c_1 Dl)}\right| + \lim_{L \to \infty} \frac{1}{L} = 0
    \]
    Por lo tanto, por el criterio de Weyl, la sucesión $(p(n))$ es equidistribuida módulo $1$ 
    cuando $c_1 \notin \mathbb{Q}$ y los demás coeficientes son racionales.

    Para el caso general, sea $p(x)=c_m x^m + c_{m-1} x^{m-1} + \dots + c_1 x + c_0$,
    sea $s = \max\{j \geq 1 : c_j \notin \mathbb{Q}\}$, es decir, $s$ es el grado
    del coeficiente irracional de mayor grado.

    Queremos demostrar que si $s\geq1$ entonces la sucesión $(p(n))$ es equidistribuida módulo $1$.
    Por~inducción sobre $s$. Para $s=1$ el resultado es cierto porque lo hemos demostrado antes.
    Supongamos que el resultado es cierto para $s-1$ y comprobemos el caso $s$.

    Para cada $h=1,2,\dots,$ definimos la función $p_h(x) = p(x+h) - p(x)$. De esta forma, 
    \[
        p_h(x) = \sum_{j=1}^m c_j ((x+h)^j - x^j)
    \]
    Desarrollando el binomio para cada $j=1,...,m$,
    \[
        (x+h)^j - x^j = \sum_{i=1}^j \binom{j}{i} x^{j-i} h^i
    \]
    Bajamos un grado a cada término, entonces
    \[
        p_h(x) = 
        \sum_{j=1}^m \sum_{i=1}^j c_j \binom{j}{i} x^{j-i} h^i =
        \sum_{j=0}^{m-1} \alpha_j x^j
    \]
    Los coeficientes $\alpha^{m-1}, \dots, \alpha^{s} \in \mathbb{Q}$, porque son combinaciones 
    lineales los términos $c_m, \dots, c_{s+1}$. Sin embargo, el coeficiente de $\alpha_{s-1}$ 
    en su combinación contiene a $c_s\notin \mathbb{Q}$, por lo que $\alpha_{s-1} \notin \mathbb{Q}$. 
    Por la hipótesis de inducción, la sucesión $(p_h(n))$ es equidistribuida módulo $1$.
    
    Por el teorema \ref{teo:diferencias_weyl}, tenemos que la sucesión $(p(n))$ es 
    equidistribuida módulo $1$ para $s\geq1$. Esto~demuestra que si existe un coeficiente irracional 
    que no sea el término independiente, entonces la sucesión $(p(n))$ es equidistribuida módulo $1$.
\end{proof}